% Options for packages loaded elsewhere
\PassOptionsToPackage{unicode}{hyperref}
\PassOptionsToPackage{hyphens}{url}
\PassOptionsToPackage{dvipsnames,svgnames,x11names}{xcolor}
%
\documentclass[
  11pt,
]{article}
\usepackage{amsmath,amssymb}
\usepackage{iftex}
\ifPDFTeX
  \usepackage[T1]{fontenc}
  \usepackage[utf8]{inputenc}
  \usepackage{textcomp} % provide euro and other symbols
\else % if luatex or xetex
  \usepackage{unicode-math} % this also loads fontspec
  \defaultfontfeatures{Scale=MatchLowercase}
  \defaultfontfeatures[\rmfamily]{Ligatures=TeX,Scale=1}
\fi
\usepackage{lmodern}
\ifPDFTeX\else
  % xetex/luatex font selection
\fi
% Use upquote if available, for straight quotes in verbatim environments
\IfFileExists{upquote.sty}{\usepackage{upquote}}{}
\IfFileExists{microtype.sty}{% use microtype if available
  \usepackage[]{microtype}
  \UseMicrotypeSet[protrusion]{basicmath} % disable protrusion for tt fonts
}{}
\makeatletter
\@ifundefined{KOMAClassName}{% if non-KOMA class
  \IfFileExists{parskip.sty}{%
    \usepackage{parskip}
  }{% else
    \setlength{\parindent}{0pt}
    \setlength{\parskip}{6pt plus 2pt minus 1pt}}
}{% if KOMA class
  \KOMAoptions{parskip=half}}
\makeatother
\usepackage{xcolor}
\usepackage[margin=2.5cm]{geometry}
\usepackage{color}
\usepackage{fancyvrb}
\newcommand{\VerbBar}{|}
\newcommand{\VERB}{\Verb[commandchars=\\\{\}]}
\DefineVerbatimEnvironment{Highlighting}{Verbatim}{commandchars=\\\{\}}
% Add ',fontsize=\small' for more characters per line
\newenvironment{Shaded}{}{}
\newcommand{\AlertTok}[1]{\textcolor[rgb]{1.00,0.00,0.00}{\textbf{#1}}}
\newcommand{\AnnotationTok}[1]{\textcolor[rgb]{0.38,0.63,0.69}{\textbf{\textit{#1}}}}
\newcommand{\AttributeTok}[1]{\textcolor[rgb]{0.49,0.56,0.16}{#1}}
\newcommand{\BaseNTok}[1]{\textcolor[rgb]{0.25,0.63,0.44}{#1}}
\newcommand{\BuiltInTok}[1]{\textcolor[rgb]{0.00,0.50,0.00}{#1}}
\newcommand{\CharTok}[1]{\textcolor[rgb]{0.25,0.44,0.63}{#1}}
\newcommand{\CommentTok}[1]{\textcolor[rgb]{0.38,0.63,0.69}{\textit{#1}}}
\newcommand{\CommentVarTok}[1]{\textcolor[rgb]{0.38,0.63,0.69}{\textbf{\textit{#1}}}}
\newcommand{\ConstantTok}[1]{\textcolor[rgb]{0.53,0.00,0.00}{#1}}
\newcommand{\ControlFlowTok}[1]{\textcolor[rgb]{0.00,0.44,0.13}{\textbf{#1}}}
\newcommand{\DataTypeTok}[1]{\textcolor[rgb]{0.56,0.13,0.00}{#1}}
\newcommand{\DecValTok}[1]{\textcolor[rgb]{0.25,0.63,0.44}{#1}}
\newcommand{\DocumentationTok}[1]{\textcolor[rgb]{0.73,0.13,0.13}{\textit{#1}}}
\newcommand{\ErrorTok}[1]{\textcolor[rgb]{1.00,0.00,0.00}{\textbf{#1}}}
\newcommand{\ExtensionTok}[1]{#1}
\newcommand{\FloatTok}[1]{\textcolor[rgb]{0.25,0.63,0.44}{#1}}
\newcommand{\FunctionTok}[1]{\textcolor[rgb]{0.02,0.16,0.49}{#1}}
\newcommand{\ImportTok}[1]{\textcolor[rgb]{0.00,0.50,0.00}{\textbf{#1}}}
\newcommand{\InformationTok}[1]{\textcolor[rgb]{0.38,0.63,0.69}{\textbf{\textit{#1}}}}
\newcommand{\KeywordTok}[1]{\textcolor[rgb]{0.00,0.44,0.13}{\textbf{#1}}}
\newcommand{\NormalTok}[1]{#1}
\newcommand{\OperatorTok}[1]{\textcolor[rgb]{0.40,0.40,0.40}{#1}}
\newcommand{\OtherTok}[1]{\textcolor[rgb]{0.00,0.44,0.13}{#1}}
\newcommand{\PreprocessorTok}[1]{\textcolor[rgb]{0.74,0.48,0.00}{#1}}
\newcommand{\RegionMarkerTok}[1]{#1}
\newcommand{\SpecialCharTok}[1]{\textcolor[rgb]{0.25,0.44,0.63}{#1}}
\newcommand{\SpecialStringTok}[1]{\textcolor[rgb]{0.73,0.40,0.53}{#1}}
\newcommand{\StringTok}[1]{\textcolor[rgb]{0.25,0.44,0.63}{#1}}
\newcommand{\VariableTok}[1]{\textcolor[rgb]{0.10,0.09,0.49}{#1}}
\newcommand{\VerbatimStringTok}[1]{\textcolor[rgb]{0.25,0.44,0.63}{#1}}
\newcommand{\WarningTok}[1]{\textcolor[rgb]{0.38,0.63,0.69}{\textbf{\textit{#1}}}}
\usepackage{longtable,booktabs,array}
\usepackage{calc} % for calculating minipage widths
% Correct order of tables after \paragraph or \subparagraph
\usepackage{etoolbox}
\makeatletter
\patchcmd\longtable{\par}{\if@noskipsec\mbox{}\fi\par}{}{}
\makeatother
% Allow footnotes in longtable head/foot
\IfFileExists{footnotehyper.sty}{\usepackage{footnotehyper}}{\usepackage{footnote}}
\makesavenoteenv{longtable}
\setlength{\emergencystretch}{3em} % prevent overfull lines
\providecommand{\tightlist}{%
  \setlength{\itemsep}{0pt}\setlength{\parskip}{0pt}}
\setcounter{secnumdepth}{5}
% Header for pandoc -> xelatex output. Liberation Serif provides bold/italic
% variants (DejaVu Serif on this system ships only Book + Bold). Unicode
% arrows / math symbols (↔, ⊗, ≤, ≥, …) that the body font lacks fall back
% to DejaVu Sans, which carries the full BMP arrow block.
\usepackage{fontspec}
\setmainfont{Liberation Serif}
\setsansfont{Liberation Sans}
\setmonofont{Liberation Mono}[Scale=MatchLowercase]

% Fallback for unicode glyphs missing in Liberation: route them through
% DejaVu Sans so we never get a tofu box in the rendered PDF.
\usepackage{newunicodechar}
\newfontfamily{\arrowfont}{DejaVu Sans}
\newunicodechar{↔}{{\arrowfont ↔}}
\newunicodechar{→}{{\arrowfont →}}
\newunicodechar{←}{{\arrowfont ←}}
\newunicodechar{⊗}{{\arrowfont ⊗}}
\newunicodechar{≤}{{\arrowfont ≤}}
\newunicodechar{≥}{{\arrowfont ≥}}
\newunicodechar{✓}{{\arrowfont ✓}}
\newunicodechar{✗}{{\arrowfont ✗}}
\newunicodechar{≈}{{\arrowfont ≈}}

\usepackage{amsmath,amssymb}
\usepackage{booktabs}
\usepackage{fancyvrb}
\ifLuaTeX
  \usepackage{selnolig}  % disable illegal ligatures
\fi
\IfFileExists{bookmark.sty}{\usepackage{bookmark}}{\usepackage{hyperref}}
\IfFileExists{xurl.sty}{\usepackage{xurl}}{} % add URL line breaks if available
\urlstyle{same}
\hypersetup{
  colorlinks=true,
  linkcolor={blue},
  filecolor={Maroon},
  citecolor={Blue},
  urlcolor={blue},
  pdfcreator={LaTeX via pandoc}}

\title{Purkinje + PVJ + Regional Ionic Numerics}
\author{cardio-dance project}
\date{2026-05}

\begin{document}
\maketitle

{
\hypersetup{linkcolor=black}
\setcounter{tocdepth}{3}
\tableofcontents
}
\hypertarget{purkinje-pvj-regional-ionic-numerics}{%
\section{Purkinje + PVJ + Regional Ionic
Numerics}\label{purkinje-pvj-regional-ionic-numerics}}

This document is the source-of-truth for the mathematical formulation
behind the Purkinje stack. Every algorithm change must update this
document in the same commit so reviewers can audit math/code drift in
one diff.

\begin{center}\rule{0.5\linewidth}{0.5pt}\end{center}

\hypertarget{domain-decomposition}{%
\subsection{1. Domain decomposition}\label{domain-decomposition}}

Three coupled discrete systems share state every PDE step:

\begin{longtable}[]{@{}
  >{\raggedright\arraybackslash}p{(\columnwidth - 4\tabcolsep) * \real{0.1889}}
  >{\raggedright\arraybackslash}p{(\columnwidth - 4\tabcolsep) * \real{0.5333}}
  >{\raggedright\arraybackslash}p{(\columnwidth - 4\tabcolsep) * \real{0.2778}}@{}}
\toprule\noalign{}
\begin{minipage}[b]{\linewidth}\raggedright
System
\end{minipage} & \begin{minipage}[b]{\linewidth}\raggedright
Discretization
\end{minipage} & \begin{minipage}[b]{\linewidth}\raggedright
Unknowns
\end{minipage} \\
\midrule\noalign{}
\endhead
\bottomrule\noalign{}
\endlastfoot
Myocardium & 3-D continuous Galerkin P1 on tetrahedra & \(V_m\) at true
DOFs \\
Purkinje cable & 1-D linear FE on a graph (subdivided edges) & \(V_p\)
at FE nodes \\
PVJ coupling & Discrete gap-junction current at terminal ↔ DOF &
Per-link current \(I_{pvj}\) \\
\end{longtable}

Time stepping is operator-split Godunov with Crank-Nicolson on diffusion
and Rush-Larsen + Forward Euler on reaction. Both PDEs march on the same
outer \(\Delta t_\mathrm{PDE}\); the cable substeps with
\(\Delta t_p \le \Delta t_\mathrm{PDE}\).

\begin{center}\rule{0.5\linewidth}{0.5pt}\end{center}

\hypertarget{sign-and-unit-conventions}{%
\subsection{2. Sign and unit
conventions}\label{sign-and-unit-conventions}}

All currents in this codebase follow the TT06 / CellML convention:

\begin{quote}
\textbf{\(I_{ion} > 0\) is outward (repolarizing).}
\textbf{\(I_{stim} < 0\) is depolarizing.}
\end{quote}

Units (consistently throughout):

\begin{longtable}[]{@{}
  >{\raggedright\arraybackslash}p{(\columnwidth - 4\tabcolsep) * \real{0.1728}}
  >{\raggedright\arraybackslash}p{(\columnwidth - 4\tabcolsep) * \real{0.2222}}
  >{\raggedright\arraybackslash}p{(\columnwidth - 4\tabcolsep) * \real{0.6049}}@{}}
\toprule\noalign{}
\begin{minipage}[b]{\linewidth}\raggedright
Quantity
\end{minipage} & \begin{minipage}[b]{\linewidth}\raggedright
Unit
\end{minipage} & \begin{minipage}[b]{\linewidth}\raggedright
Notes
\end{minipage} \\
\midrule\noalign{}
\endhead
\bottomrule\noalign{}
\endlastfoot
Length & mm & mesh coordinates, fiber lengths \\
Time & ms & \\
Voltage & mV & Cell membrane potential \\
Concentration & mM (mmol/L) & \([Na^+]_i\), \([Ca^{2+}]_{SS}\), etc. \\
\(C_m\) & μF/mm² & membrane capacitance per unit area \\
\(\chi\) & 1/mm & surface-to-volume ratio \\
\(\sigma\) & mS/mm & tissue conductivity \\
\(g_\mathrm{ion}\) & mS/μF & normalized cell conductance \\
\(I_{ion}, I_{stim}\) & μA/μF & normalized membrane current (∝ dV/dt) \\
\end{longtable}

\textbf{Verification:} the monodomain RHS in
\texttt{MonodomainStepper::BuildRhs} reads
\texttt{-chi*Cm*M*(I\_ion\ +\ I\_stim\ +\ I\_pvj)} --- the negative sign
and \texttt{chi*Cm} prefactor convert from ``current per unit area''
semantics to units consistent with the Crank-Nicolson system
\(A V^{n+1} = B V^n + \mathrm{RHS}\).

\begin{center}\rule{0.5\linewidth}{0.5pt}\end{center}

\hypertarget{myocardial-monodomain-existing-recapped-here}{%
\subsection{3. Myocardial monodomain (existing, recapped
here)}\label{myocardial-monodomain-existing-recapped-here}}

Continuous form on heart domain \(\Omega_h\):
\[\chi C_m \partial_t V_m = \nabla\cdot(\boldsymbol\sigma \nabla V_m) - \chi C_m (I_{ion}(V_m, \mathbf{s}) + I_{stim} + I_{pvj}^h),\]
\[\partial_t \mathbf{s} = \mathbf{f}(V_m, \mathbf{s}).\]

Spatial discretization: P1 CG, mass \(M\) and stiffness \(K\) assembled
with the fiber-aware tensor
\(\boldsymbol\sigma = \sigma_f\, \mathbf{f}\!\otimes\!\mathbf{f} + \sigma_s\, \mathbf{s}\!\otimes\!\mathbf{s} + \sigma_n\, \mathbf{n}\!\otimes\!\mathbf{n}\).
With \textbf{AV-delay / fibrosis regions} the stiffness uses an
attribute-wise multiplier \(\alpha(x)\):
\[K_{ij} = \int_{\Omega_h} \alpha(x)\, (\nabla\phi_i)^T \boldsymbol\sigma \nabla\phi_j\,\mathrm{d}x,\]
where \$\alpha = \$ \texttt{av\_delay\_sigma\_scale} (0.02) on AV-delay
attributes, \texttt{fibrosis\_sigma\_scale} (0.1) on fibrosis
attributes, 1.0 elsewhere.

\textbf{Code mapping:} \texttt{Assembler.cpp} builds
\texttt{region\_sigma\_scale\_} as a \texttt{PWConstCoefficient} indexed
by mesh attribute, then wraps the \texttt{FiberTensorCoefficient} in
\texttt{ScalarMatrixProductCoefficient}. Single DiffusionIntegrator
consumes the product.

\begin{center}\rule{0.5\linewidth}{0.5pt}\end{center}

\hypertarget{purkinje-cable}{%
\subsection{4. Purkinje cable}\label{purkinje-cable}}

\hypertarget{continuous-form}{%
\subsubsection{4.1 Continuous form}\label{continuous-form}}

Each Purkinje fiber is modeled as a 1-D cable. With per-edge axial
conductivity \(g_a\) (mS/mm) and uniform \(C_m^p\) (μF/mm²),
\[C_m^p \partial_t V_p = \partial_x (g_a \partial_x V_p) - I_{ion}^p(V_p, \mathbf{s}_p) - I_{pvj}^p - I_{stim}^p.\]
At branch points (graph vertices with degree \textgreater{} 2) currents
are conserved by construction: the FE assembly's nodal balance enforces
Kirchhoff's law because each incident edge contributes its own
\(g_a/L_\mathrm{seg}\) entry to the same row.

\hypertarget{spatial-discretization}{%
\subsubsection{4.2 Spatial
discretization}\label{spatial-discretization}}

Each graph edge of length \(L\) is subdivided into \(k\) FE sub-segments
of length \(L/k\). Linear FE on each sub-segment yield the local
stiffness
\[K^e_{\mathrm{loc}} = \frac{g_a}{L/k} \begin{pmatrix} 1 & -1 \\ -1 & 1 \end{pmatrix}\]
and a lumped mass \(M^e_{\mathrm{loc}} = \tfrac{C_m^p L/k}{2} I\).
Global \(K_p\) and \(\mathrm{diag}(M_p)\) are sparse-assembled in
\texttt{PurkinjeCableSolver::BuildSparseMatrices}.

\textbf{Edge weight semantics:} the third token in \texttt{.network}
edge lines is the \textbf{axial conductance per mm of edge length},
i.e.~\(g_a\). If absent, the value defaults to
\texttt{cfg.purkinje\_edge\_g\_mS\_per\_mm}. The per-segment conductance
contribution is then \(g_\mathrm{seg} = g_a / L_\mathrm{seg}\).

\hypertarget{time-stepping}{%
\subsubsection{4.3 Time stepping}\label{time-stepping}}

Operator split per outer \(\Delta t_\mathrm{PDE}\):

\begin{enumerate}
\def\labelenumi{\arabic{enumi}.}
\tightlist
\item
  Subdivide into
  \(n_\mathrm{sub} = \lceil \Delta t_\mathrm{PDE} / \Delta t_p\rceil\)
  steps of \(\Delta t = \Delta t_\mathrm{PDE}/n_\mathrm{sub}\).
\item
  Per substep:

  \begin{enumerate}
  \def\labelenumii{\arabic{enumii}.}
  \tightlist
  \item
    Compute \(I_{ion}^p(V_p, \mathbf{s}_p)\) at current \(V_p\).
  \item
    Advance \(\mathbf{s}_p\) via Rush-Larsen (gates) + Forward Euler
    (concentrations) using the current \(V_p\).
  \item
    Build external current \(I_\mathrm{ext} = I_{pvj}^p + I_{stim}^p\).
  \item
    Solve
    \((M_p + \tfrac{\Delta t}{2} K_p) V_p^{n+1} = (M_p - \tfrac{\Delta t}{2} K_p) V_p^n - \Delta t\, M_p (I_{ion}^p + I_\mathrm{ext})\).
  \end{enumerate}
\end{enumerate}

Linear solve: CG with Jacobi (\texttt{mfem::DSmoother}) preconditioner;
the system matrices are cached and rebuilt only when \(\Delta t\)
changes.

\hypertarget{stability-dt-selection}{%
\subsubsection{4.4 Stability / dt
selection}\label{stability-dt-selection}}

Crank-Nicolson is unconditionally A-stable for the diffusion split.
Reaction substep is governed by the fastest gating \(\tau\) (I\_Na:
\(\tau_m \approx 0.05\) ms near upstroke). Practical guidance: -
\(\Delta t_p \le 0.01\) ms during upstroke (handled by Rush-Larsen exact
gate update; FE on concentrations less sensitive) -
\(\Delta t_\mathrm{PDE}/\Delta t_p \le 4\) keeps PVJ delay quantization
\textless{} 0.005 ms

\begin{center}\rule{0.5\linewidth}{0.5pt}\end{center}

\hypertarget{purkinje-ionic-model-stewart-2009}{%
\subsection{5. Purkinje ionic model (Stewart
2009)}\label{purkinje-ionic-model-stewart-2009}}

\hypertarget{reference}{%
\subsubsection{5.1 Reference}\label{reference}}

Stewart, P., Aslanidi, O. V., Noble, D., Noble, P. J., Boyett, M. R., \&
Zhang, H. (2009). \emph{Mathematical models of the electrical action
potential of Purkinje fibre cells.} Phil. Trans. R. Soc. A 367 (1896),
2225--2255. DOI 10.1098/rsta.2008.0283.

The Stewart model is itself a Purkinje variant of TT06 with HCN funny
current \(I_f\) added and modified \(I_{to}, I_{Ks}, I_{K1}\) kinetics.

\hypertarget{currents-target-list}{%
\subsubsection{5.2 Currents (target list)}\label{currents-target-list}}

\(I_\mathrm{ion} = I_{Na} + I_{CaL} + I_{to} + I_{Kr} + I_{Ks} + I_{K1} + I_f^{Na} + I_f^K + I_{NaCa} + I_{NaK} + I_{pCa} + I_{pK} + I_{bNa} + I_{bCa} + I_\mathrm{sus}\)

\textbf{I\_CaL (TT06/Stewart shifted GHK):}
\[I_{CaL} = g_{CaL}\, d f f_2 f_{Cass}\, \frac{4(V-15)F^2}{RT}\, \frac{0.25\,[Ca]_{SS}\,e^{2(V-15)F/RT} - [Ca]_o}{e^{2(V-15)F/RT} - 1}.\]

\textbf{Important:} both the prefactor and the exponential argument use
the shifted voltage \(V_\mathrm{shift} = V - 15\) mV. A historical bug
in this codebase mixed \(V-15\) in the prefactor with \(V\) in the
exponent and produced huge non-physiological CaL drive.

\hypertarget{state-vector-and-integration}{%
\subsubsection{5.3 State vector and
integration}\label{state-vector-and-integration}}

20 states:
\texttt{{[}0{]}V,\ {[}1{]}m,\ {[}2{]}h,\ {[}3{]}j,\ {[}4{]}xr1,\ {[}5{]}xr2,\ {[}6{]}xs,\ {[}7{]}r,\ {[}8{]}s,\ {[}9{]}d,\ {[}10{]}f,\ {[}11{]}f2,\ {[}12{]}fCass,\ {[}13{]}y(I\_f),\ {[}14{]}Cai,\ {[}15{]}CaSR,\ {[}16{]}CaSS,\ {[}17{]}Nai,\ {[}18{]}Ki,\ {[}19{]}Rprime}.

Per substep: - States 1..13 (gates): Rush-Larsen
\(x \leftarrow x_\infty - (x_\infty - x)\,e^{-\Delta t/\tau}\). - States
14..19 (concentrations + RyR): Forward Euler
\(x \leftarrow x + \Delta t\,\dot x\), with positivity clamps on
\([Ca^{2+}]_*, [Na^+]_i, [K^+]_i\).

\hypertarget{status-2026-05-cellml-codegen-verified}{%
\subsubsection{5.4 Status (2026-05): CellML codegen
verified}\label{status-2026-05-cellml-codegen-verified}}

\texttt{StewartPurkinjeModel} now delegates to the \textbf{official
CellML codegen} in \texttt{src/ode/stewart\_generated.\{h,c\}}, lifted
verbatim from
\texttt{models.cellml.org/exposure/38cf8387b0707f0ef6947f009710aeb5/\ stewart\_aslanidi\_noble\_noble\_boyett\_zhang\_2009.cellml/@@cellml\_codegen/C}.
Functions are prefixed \texttt{stewart\_} to avoid linker collision with
TT06's identical signatures.

\hypertarget{single-cell-verification-test_stewart_single_cell}{%
\paragraph{\texorpdfstring{Single-cell verification
(\texttt{test\_stewart\_single\_cell})}{Single-cell verification (test\_stewart\_single\_cell)}}\label{single-cell-verification-test_stewart_single_cell}}

Brief (-52 μA/μF × 0.5 ms) stim from MDP=-90 mV, free run for 1000 ms.
Codegen output:

\begin{longtable}[]{@{}
  >{\raggedright\arraybackslash}p{(\columnwidth - 6\tabcolsep) * \real{0.2593}}
  >{\raggedright\arraybackslash}p{(\columnwidth - 6\tabcolsep) * \real{0.2593}}
  >{\raggedright\arraybackslash}p{(\columnwidth - 6\tabcolsep) * \real{0.2593}}
  >{\raggedright\arraybackslash}p{(\columnwidth - 6\tabcolsep) * \real{0.2222}}@{}}
\toprule\noalign{}
\begin{minipage}[b]{\linewidth}\raggedright
Quantity
\end{minipage} & \begin{minipage}[b]{\linewidth}\raggedright
Paper Fig. 2
\end{minipage} & \begin{minipage}[b]{\linewidth}\raggedright
Current code
\end{minipage} & \begin{minipage}[b]{\linewidth}\raggedright
Test bound
\end{minipage} \\
\midrule\noalign{}
\endhead
\bottomrule\noalign{}
\endlastfoot
V\_rest\_pre & -91 mV (paced) & -74.2 (HCN-driven natural rest) &
{[}-78, -68{]} ✓ \\
V\_peak & +30 mV (plateau) & +53.7 (single-stim spike) & {[}10, 60{]}
✓ \\
V\_plateau & +30 mV & +17.2 (10 ms post-stim) & {[}10, 40{]} ✓ \\
APD90 & 340 ms & 293 ms & {[}200, 400{]} ✓ \\
V\_at\_min & -91 mV (paced) & -75.9 (post-AP rest) & \textless{} -70
✓ \\
\end{longtable}

The ``natural'' Stewart rest is \textbf{-74 mV}, not -91 mV: HCN/I\_f
drive keeps the cell perpetually rising slowly toward firing threshold
(this is expected Purkinje autorhythmic behavior). The paper's -91 mV
figures come from steady-state pacing where the cell has been clamped
low between beats; a single-stim test from MDP=-90 instead settles to
V\_rest=-74 mV. We test against the codegen's actual single-stim
morphology, not the steady-state pacing values.

\hypertarget{cable-verification-test_purkinje_cable_cv}{%
\paragraph{\texorpdfstring{Cable verification
(\texttt{test\_purkinje\_cable\_cv})}{Cable verification (test\_purkinje\_cable\_cv)}}\label{cable-verification-test_purkinje_cable_cv}}

51-node Stewart cable, 1 mm spacing (50 mm total), Crank-Nicolson +
Rush- Larsen integration with
\texttt{cfg.purkinje\_edge\_g\_mS\_per\_mm\ =\ 0.05} (lumped edge axial
conductance tuned for L\_seg=1 mm, C\_m=0.01 μF/mm² to give CV in the
published Purkinje range). Cluster-stim of nodes 0,1,2 at -150 μA/μF for
1 ms; activation measured by V \textgreater{} -40 mV at probe nodes 10
and 40 (30 mm separation).

\begin{longtable}[]{@{}llll@{}}
\toprule\noalign{}
Quantity & Published Purkinje & Current code & Test bound \\
\midrule\noalign{}
\endhead
\bottomrule\noalign{}
\endlastfoot
CV & 1.5--4 m/s & 5.4 m/s & {[}1.5, 6.5{]} ✓ \\
\end{longtable}

CV depends sensitively on the lumped edge conductance \texttt{g\_a} and
on the discretization spacing; the test's 5.4 m/s is at the high end of
published values but well inside the physiologically plausible window.
For a tighter target, lower \texttt{purkinje\_edge\_g\_mS\_per\_mm}
further (CV \textasciitilde{} sqrt(g\_a/C\_m)).

\hypertarget{diagnostic-tooling}{%
\paragraph{Diagnostic tooling}\label{diagnostic-tooling}}

\texttt{tools/debug\_stewart\_currents} now drives the production model
directly (no shadow-integration); it dumps \texttt{t,V,I\_total} CSV at
0.05 ms intervals. Used to verify the codegen-driven AP shape visually
after any future refactor that touches Stewart integration.

\hypertarget{importing-cellml-reference-code}{%
\subsubsection{5.5 Importing CellML reference
code}\label{importing-cellml-reference-code}}

Long-term remediation: 1. Download Stewart 2009 from
\texttt{https://models.physiomeproject.org/exposure/...} (CellML 1.0
XML). 2. Run OpenCOR or \texttt{cellml-tools} to emit C with
\texttt{initConsts}, \texttt{computeRates}, \texttt{computeVariables}
matching the TT06 generated code style. 3. Replace
\texttt{src/ode/StewartPurkinjeModel.cpp} body with thin wrappers
delegating to the generated functions, mirroring the
\texttt{tt06\_generated.\{h,c\}} pattern. Keep the gate/concentration
index lists (\texttt{kGateMaps}) as the only hand-written part. 4.
Re-enable the strict physiological asserts in
\texttt{test\_stewart\_single\_cell.cpp} and remove the
\texttt{WILL\_FAIL} marker.

\begin{center}\rule{0.5\linewidth}{0.5pt}\end{center}

\hypertarget{atrial-ionic-model-grandi-2011}{%
\subsection{6. Atrial ionic model (Grandi
2011)}\label{atrial-ionic-model-grandi-2011}}

\texttt{Grandi2011Model} is a \textbf{reduced 15-state compact port}
retaining the dominant atrial currents:
\(I_{Na}, I_{CaL}^\mathrm{simplified}, I_{Kr}, I_{Ks}, I_{to,f}, I_{Kur}, I_{K1}, I_{NaK}, I_{NaCa}, I_{pCa}, I_{bNa}, I_{bCa}\).
The simplified \(I_{CaL}\) uses a constant driving force \((V - 60)\)
rather than full GHK, matching what is typical for ``minimal-model''
atrial sims when only conduction (not Ca dynamics) matters.

\textbf{Status (2026-05):} no single-cell verification test exists yet
for Grandi. Use this model only for atrial-region conduction studies,
not for APD/AERP analysis.

\textbf{Reference (target):} Grandi, E., Pandit, S. V., Voigt, N.,
Workman, A. J., Dobrev, D., Jalife, J., \& Bers, D. M. (2011).
\emph{Human atrial action potential and Ca²⁺ model: sinus rhythm and
chronic atrial fibrillation.} Circulation Research 109 (9), 1055--1066.

\begin{center}\rule{0.5\linewidth}{0.5pt}\end{center}

\hypertarget{passive-leak-model}{%
\subsection{7. Passive (leak) model}\label{passive-leak-model}}

\[I_{ion}(V) = g_\mathrm{leak}\, (V - V_\mathrm{rest}).\]

Used in AV-delay and fibrosis regions to suppress active depolarization
while allowing electrotonic spread. State-free, so checkpoint is just
the two parameters \((g_\mathrm{leak}, V_\mathrm{rest})\).

\begin{center}\rule{0.5\linewidth}{0.5pt}\end{center}

\hypertarget{pvj-coupling}{%
\subsection{8. PVJ coupling}\label{pvj-coupling}}

\hypertarget{continuous-form-1}{%
\subsubsection{8.1 Continuous form}\label{continuous-form-1}}

For each terminal Purkinje node \(t\) mapped to nearest myocardial true
DOF \(d(t)\): \[I_{pvj,t} = g_{pvj}\, [V_p^t - V_m^{d(t)}]\, s_{pvj}.\]
- \(g_{pvj}\) in mS/μF (config: \texttt{pvj\_g\_mS}, default 0.5) -
\(s_{pvj}\) unitless scale (config: \texttt{pvj\_current\_scale},
default 1.0)

\hypertarget{sign-convention}{%
\subsubsection{8.2 Sign convention}\label{sign-convention}}

\begin{quote}
When \(V_p > V_m\) (Purkinje more depolarized), \(I_{pvj} > 0\). The
\textbf{myocardial RHS contribution} is \(+I_{pvj}\) at DOF \(d(t)\),
treated with the same \texttt{+} sign as \(I_{ion}\) (outward
semantics). So the actual contribution to the RHS is
\(-\chi C_m M (\ldots + I_{pvj})\), i.e.~when \(V_p > V_m\) the negative
sign \textbf{depolarizes} \(V_m\). ✓
\end{quote}

\begin{quote}
The \textbf{Purkinje cable} sees \(-I_{pvj}\) at terminal \(t\)
(Newton's third law: charge flows from Purkinje to myocardium, so
Purkinje loses positive ion flux). Implementation in
\texttt{PurkinjeCableSolver::Advance}:
\texttt{i\_pvj\_p{[}fe{]}\ =\ pvj\_g\ *\ scale\ *\ (vp{[}fe{]}\ -\ Vm)}.
Same sign as \texttt{BuildHeartCouplingCurrent} because both are added
to their respective \texttt{I\_ion} term.
\end{quote}

\hypertarget{mapping-algorithm}{%
\subsubsection{8.3 Mapping algorithm}\label{mapping-algorithm}}

\texttt{PvjCoupler::BuildMapping}: 1. Each rank computes nearest local
true DOF for each terminal. 2.
\texttt{MPI\_Allreduce(MPI\_DOUBLE\_INT,\ MPI\_MINLOC)} selects the
global owner. 3. Terminal-DOF pairs farther than
\texttt{pvj\_max\_dist\_mm} are dropped. 4. Each owner appends a
\texttt{LocalPvjLink\ \{purkinje\_node,\ heart\_tdof,\ dist\}} to its
\texttt{local\_links\_}.

Coordinate extraction uses MFEM \texttt{FunctionCoefficient} projected
onto the P1 space, then \texttt{GetTrueDofs} --- guarantees exact match
with monodomain DOFs.

\hypertarget{optional-anatomical-delay}{%
\subsubsection{8.4 Optional anatomical
delay}\label{optional-anatomical-delay}}

When \texttt{pvj\_delay\_ms\ \textgreater{}\ 0}: - After each Purkinje
step, the current \(V_m\) at terminals is appended to a ring buffer
(\texttt{PvjCoupler::AppendDelayedSample}). - The cable advance reads
the sample with timestamp \(\le t_\mathrm{mid} -  \mathrm{delay}\)
(\texttt{ReadDelayedSample}) --- implements a transport delay of
\texttt{delay\_ms} between heart and Purkinje sampling.

Buffer trim keeps history within \(4\times\) delay for safety.

\hypertarget{smeared-injection-3d-source-sink-mismatch-fix}{%
\subsubsection{8.5 Smeared injection (3D source-sink mismatch
fix)}\label{smeared-injection-3d-source-sink-mismatch-fix}}

Under default single-DOF injection, the heart-side current is added to
one true DOF nearest the terminal. On 3D unstructured tet meshes with
cell size \(h\) comparable to or larger than the cardiac space constant
\(\lambda = \sqrt{D \tau_\mathrm{AP}} \approx 0.3\) mm, this point
source is drained by 8-12 surrounding resting DOFs \emph{before} I\_Na
can establish locally --- the activated patch is too small to
self-sustain. Symptom: \(V_m\) briefly peaks near the injection then
collapses back without propagation.

\texttt{pvj\_smear\_radius\_mm\ \textgreater{}\ 0} distributes the
per-terminal current over a ball of radius \(R\) around the anchor DOF:

\[I_{pvj}^{(i)} = w_i\, g_{pvj}\, s_{pvj}\, (V_m^{(i)} - V_p)
\quad\text{for all } i \in \mathcal{B}_R(\mathrm{anchor})\]

with \(w_i = 1 / |\mathcal{B}_R|\) uniform; the per-terminal total
injection magnitude is unchanged. Anchor DOF is still used (single owner
rank) for \(V_m\) feedback to the cable, so the MPI\_Allreduce SUM in
\texttt{AdvancePurkinje} still recovers exactly one owner-side value per
terminal even though the smeared list contains multiple links.

Tested on the half-ellipsoid demo: with
\texttt{pvj\_smear\_radius\_mm=0} the wave collapses, with
\texttt{pvj\_smear\_radius\_mm=4} (full-scale) or \texttt{=2}
(half-scale) it propagates outward through the wall.

\begin{center}\rule{0.5\linewidth}{0.5pt}\end{center}

\hypertarget{regional-ionic-dispatcher}{%
\subsection{9. Regional ionic
dispatcher}\label{regional-ionic-dispatcher}}

\hypertarget{mapping}{%
\subsubsection{9.1 Mapping}\label{mapping}}

Mesh element attributes drive region assignment:

\begin{longtable}[]{@{}lll@{}}
\toprule\noalign{}
Region & Cell model & Default config tag \\
\midrule\noalign{}
\endhead
\bottomrule\noalign{}
\endlastfoot
Ventricle & TT06 & \texttt{ventricle\_volume\_attrs} \\
Atria & Grandi 2011 & \texttt{atria\_volume\_attrs} \\
AV-delay & Passive & \texttt{av\_delay\_volume\_attrs} \\
Fibrosis & Passive & \texttt{fibrosis\_volume\_attrs} \\
\end{longtable}

For each owned true DOF, the region with the \textbf{highest priority}
among incident elements wins (priority order: AV-delay \textgreater{}
Fibrosis \textgreater{} Atria \textgreater{} Ventricle). This guards
against a DOF on a region boundary getting an inappropriate active
model.

\hypertarget{state-isolation-resolved-2026-05}{%
\subsubsection{9.2 State isolation (resolved
2026-05)}\label{state-isolation-resolved-2026-05}}

Each child model is constructed with size =
\texttt{LocalDofCount(region)}, NOT the full DOF count.
\texttt{RegionalIonicModel::ComputeIion} and \texttt{AdvanceStates}
gather \(V_m\) into a region-local buffer using
\texttt{region\_indices\_{[}r{]}}, invoke the child on that buffer, and
scatter the result back into the global I\_ion vector. So a child's
per-DOF state index is \textbf{region-local}, never indexed by global
true DOF.

Empty regions still allocate a 1-DOF placeholder so the unique\_ptr is
constructible, but that placeholder is never invoked because
\texttt{per\_region\_count\_{[}r{]}\ ==\ 0} short-circuits the run path.

Tested in \texttt{tests/test\_regional\_isolation.cpp}: a 1-D 4-element
mesh with ventricle on elements 0-1 (TT06) and AV-delay on elements 2-3
(Passive) verifies that (a) DOFs partition the tdof set, (b) Passive at
V\_rest yields exactly 0 I\_ion, (c) repeated AdvanceStates with extreme
ventricular V does not corrupt AV-delay output.

\begin{center}\rule{0.5\linewidth}{0.5pt}\end{center}

\hypertarget{verification-matrix}{%
\subsection{10. Verification matrix}\label{verification-matrix}}

Required tests (see \texttt{tests/}):

\begin{longtable}[]{@{}
  >{\raggedright\arraybackslash}p{(\columnwidth - 4\tabcolsep) * \real{0.3977}}
  >{\raggedright\arraybackslash}p{(\columnwidth - 4\tabcolsep) * \real{0.4886}}
  >{\raggedright\arraybackslash}p{(\columnwidth - 4\tabcolsep) * \real{0.1136}}@{}}
\toprule\noalign{}
\begin{minipage}[b]{\linewidth}\raggedright
Test
\end{minipage} & \begin{minipage}[b]{\linewidth}\raggedright
Scope
\end{minipage} & \begin{minipage}[b]{\linewidth}\raggedright
Status
\end{minipage} \\
\midrule\noalign{}
\endhead
\bottomrule\noalign{}
\endlastfoot
\texttt{test\_tt06\_single\_cell} & TT06 AP shape (existing) & ✓ pass \\
\texttt{test\_stewart\_single\_cell} & Stewart codegen AP morphology & ✓
pass \\
\texttt{test\_purkinje\_cable\_smoke} & Cable diffusion correctness
(passive) & ✓ pass \\
\texttt{test\_purkinje\_cable\_cv} & 50 mm Stewart cable CV in {[}1.5,
6.5{]} m/s & ✓ pass \\
\texttt{test\_pvj\_active\_sign} & PVJ depolarizes myocardium when
\(V_p>V_m\) & ✓ pass \\
\texttt{test\_regional\_isolation} & Regional dispatch state isolation &
✓ pass \\
\texttt{test\_checkpoint\_model\_id} & CheckpointIO refuses cross-model
load & ✓ pass \\
\texttt{test\_grandi2011\_single\_cell} & Atrial AP shape & TODO \\
\end{longtable}

End-to-end: - Full Niederer benchmark with regional + Purkinje (pending
mesh tools). - Calibration sweep over \(g_{pvj}, \sigma_\mathrm{AV}\) vs
target activation times (script: \texttt{tools/calibrate\_purkinje.py},
TODO).

\begin{center}\rule{0.5\linewidth}{0.5pt}\end{center}

\hypertarget{checkpoint-format}{%
\subsection{11. Checkpoint format}\label{checkpoint-format}}

\hypertarget{meta-file-chkpt_v3}{%
\subsubsection{11.1 Meta file (CHKPT\_V3)}\label{meta-file-chkpt_v3}}

Plain-text, written only by rank 0:

\begin{verbatim}
CHKPT_V3
<world_size>
<step>
<t_ms>
<model_id>
\end{verbatim}

\texttt{model\_id} is \texttt{IIonicModel::ModelId()} of the model that
wrote the checkpoint. \texttt{LoadLatest} refuses to proceed if the
runtime model's \texttt{ModelId()} does not match (or, for legacy V2/V0
metas without \texttt{model\_id}, if the runtime model is not TT06).

\hypertarget{shards}{%
\subsubsection{11.2 Shards}\label{shards}}

\begin{itemize}
\tightlist
\item
  \texttt{vm\_rank\textless{}NNNNNN\textgreater{}.gf} ---
  \texttt{mfem::ParGridFunction::Save} per rank.
\item
  \texttt{cell\_rank\textless{}NNNNNN\textgreater{}.bin} ---
  \texttt{IIonicModel::SaveState} per rank. Binary layout is
  model-specific; cross-model loads are caught by the model\_id check
  before any binary is read.
\end{itemize}

\hypertarget{format-compatibility}{%
\subsubsection{11.3 Format compatibility}\label{format-compatibility}}

\begin{longtable}[]{@{}ll@{}}
\toprule\noalign{}
Meta header & Status \\
\midrule\noalign{}
\endhead
\bottomrule\noalign{}
\endlastfoot
CHKPT\_V3 & current, supported \\
CHKPT\_V2 & TT06-only legacy; refused otherwise \\
(legacy) & single-rank TT06 only \\
\end{longtable}

\hypertarget{verified-by}{%
\subsubsection{11.4 Verified by}\label{verified-by}}

\texttt{tests/test\_checkpoint\_model\_id.cpp}: writes a TT06
checkpoint, then asserts \texttt{PassiveModel} LoadLatest returns false
and a fresh \texttt{TT06Model} LoadLatest returns true with matching
step/t\_ms.

\begin{center}\rule{0.5\linewidth}{0.5pt}\end{center}

\hypertarget{change-log-discipline}{%
\subsection{12. Change-log discipline}\label{change-log-discipline}}

Every commit touching algorithms in this stack must include a one-line
note in the message body:

\begin{verbatim}
Math:    <equation/property changed>
Code:    <file/function changed>
Mapping: <how the verification stays sound>
\end{verbatim}

If a verification test no longer passes after a change, mark it
\texttt{WILL\_FAIL} in the same commit and add a TODO line; do not
silently relax the assertion.

\begin{center}\rule{0.5\linewidth}{0.5pt}\end{center}

\hypertarget{fiber-field-current-limitation}{%
\subsection{13. Fiber field (current
limitation)}\label{fiber-field-current-limitation}}

The fiber tensor coefficient
\(\boldsymbol\sigma = \sigma_f\, \mathbf{f}\!\otimes\!\mathbf{f} + \sigma_s\, \mathbf{s}\!\otimes\!\mathbf{s} + \sigma_n\, \mathbf{n}\!\otimes\!\mathbf{n}\)
is fully wired (\texttt{Assembler::InitializeFiberCoefficients} reads
three \texttt{.gf} files when \texttt{use\_fiber\_gf=1}). Conduction
velocity along fiber vs transverse scales as
\(\sqrt{\sigma_f / \sigma_t} = \sqrt{0.1334/0.0176} \approx 2.75\times\)
for the project default values.

\textbf{However} the \texttt{.gf} files we currently emit (via
\texttt{tools/finalize\_half\_ellipsoid\_mesh.cpp\ -\/-emit-fibers}) are
\textbf{constant vectors} \(\mathbf{f}=(1,0,0)\),
\(\mathbf{s}=(0,1,0)\), \(\mathbf{n}=(0,0,1)\) at every DOF. The
conductivity tensor is therefore fixed Cartesian diagonal
\(\mathrm{diag}(\sigma_f, \sigma_s, \sigma_n)\) --- anisotropic along
\(\hat{x}\) but \textbf{not} following the natural Streeter helix of
real left ventricles (\(-60°\) at endo rotating to \(+60°\) at epi).

What this is sufficient for: - PVJ sign / magnitude verification - PetsC
ASM iteration-count benchmark - Pipeline correctness demo (mesh +
Purkinje + ECG) - Order-of-magnitude QRS duration

What this is \textbf{not} sufficient for: - Realistic T-wave morphology
(lacks transmural repolarization gradient driven by helical fibers) -
Multi-lead ECG differences - Bundle-branch-block clinical case studies

\textbf{Adding rule-based fiber generation (TODO)}: implement
Bayer-Blake- Trayanova 2012 LDRBM via four Laplace-Dirichlet solves
(\(\phi_\mathrm{epi}, \phi_\mathrm{endo}, \phi_\mathrm{base}, \phi_\mathrm{long}\)),
construct local transmural / longitudinal / circumferential coordinate
frames from \(\nabla\phi\), then rotate the sheet vector by the Streeter
helix angle
\(\alpha(d/W) = \alpha_\mathrm{endo} + (d/W)\, (\alpha_\mathrm{epi} - \alpha_\mathrm{endo})\)
across wall fraction \(d/W\). \textasciitilde200 lines. Not implemented.

\begin{center}\rule{0.5\linewidth}{0.5pt}\end{center}

\hypertarget{half-ellipsoid-lv-demo}{%
\subsection{14. Half-ellipsoid LV demo}\label{half-ellipsoid-lv-demo}}

A worked end-to-end pipeline lives in
\texttt{config/half\_ellipsoid\_purkinje.options} and produces the
figures / movies / pseudo-ECG under
\texttt{output/half\_ellipsoid\_purkinje/}.

\hypertarget{geometry}{%
\subsubsection{14.1 Geometry}\label{geometry}}

Half-ellipsoid hollow LV (open at base \(x=0\)): - Outer ellipsoid:
\(a=35, b=22, c=22\) mm - Inner cavity:
\(a_\mathrm{in}=28, b_\mathrm{in}=15, c_\mathrm{in}=15\) mm - Wall
thickness ranges \(\sim 7\) mm equator, tapering at apex.

Mesh by \texttt{tools/half\_ellipsoid.geo} (gmsh OpenCASCADE CSG →
Frontal- Delaunay 2D + Delaunay 3D). At \(h=1.0\) mm: 21,707 nodes /
124,715 tets. \texttt{tools/finalize\_half\_ellipsoid\_mesh.cpp}
reclassifies boundary triangles into epi / endo / base by true
face-centroid distance.

A previous Cartesian-carve mesh tool
(\texttt{generate\_half\_ellipsoid\_case}) exists for reference but
produces stair-stepped boundaries; superseded.

\hypertarget{purkinje-network}{%
\subsubsection{14.2 Purkinje network}\label{purkinje-network}}

\texttt{tools/generate\_purkinje\_tree.py} (Costabal-style fractal):
from a root near the apex along \(-\hat{x}\), with 30° bifurcation
angles, depth 9, bifurcation probability 0.85. Default at full scale:
192 graph nodes / 91 terminals. The bbox argument confines growth to
within the cavity.

\hypertarget{pvj-smear}{%
\subsubsection{14.3 PVJ + smear}\label{pvj-smear}}

Activation comes from Purkinje only (no direct myocardial stim). Cluster
pace at root nodes 0/1/2 (-150 µA/µF for 1 ms) drives the cable; PVJ
maps each terminal to the nearest endocardial DOF (44/91 typically map
within \texttt{pvj\_max\_dist\_mm=5}). Heart-side injection is smeared
over a 4 mm ball (full scale) or 2 mm (half scale) to clear the
source-sink hurdle.

\hypertarget{pseudo-ecg-probe}{%
\subsubsection{14.4 Pseudo-ECG probe}\label{pseudo-ecg-probe}}

\texttt{PseudoEcg} evaluates
\(\phi(\vec x_p, t) = \frac{\sigma_i}{4\pi\sigma_b} \sum_e V_e\, \nabla V_m^{(e)} \cdot (\vec x_p - \vec c_e) / |\vec x_p - \vec c_e|^3\)
at element centroids, MPI\_Allreduce SUM, write one CSV column per
probe. Default config places one lead at \((17.5, 80, 0)\) mm.

\hypertarget{run-results-single-rank-plain-cg-full-scale-h1mm}{%
\subsubsection{14.5 Run results (single rank, plain CG; full scale
h=1mm)}\label{run-results-single-rank-plain-cg-full-scale-h1mm}}

\begin{longtable}[]{@{}ll@{}}
\toprule\noalign{}
Quantity at \(t=100\) ms & Value \\
\midrule\noalign{}
\endhead
\bottomrule\noalign{}
\endlastfoot
Wall (100 ms simulated) & 5m52s \\
KSP iter / step (plain CG) & 50 ± 6 \\
KSP iter / step (PETSc ASM) & 5 ± 1 \\
\(V_m\) min / max / mean & -93.7 / +31.3 / -66.6 mV \\
Purkinje terminals mapped & 44 / 91 \\
\end{longtable}

Visualisation artifacts (committed to git for inspection): -
\texttt{output/half\_ellipsoid\_purkinje/vm\_movie.mp4} -
\texttt{output/half\_ellipsoid\_purkinje/vm\_movie\_clipped.mp4} -
\texttt{output/half\_ellipsoid\_purkinje/pseudo\_ecg.png} -
\texttt{output/half\_ellipsoid\_purkinje/gallery/*.png} (mesh wireframe
+ cross-sections)

\begin{center}\rule{0.5\linewidth}{0.5pt}\end{center}

\hypertarget{linear-solver-crank-nicolson-diffusion-step}{%
\subsection{15. Linear solver (Crank-Nicolson diffusion
step)}\label{linear-solver-crank-nicolson-diffusion-step}}

The CN system \(A V^{n+1} = \mathrm{RHS}\) with
\(A = \chi C_m / \Delta t \cdot M + \tfrac{1}{2} K\) is symmetric
positive- definite (M and K both SPD; positive linear combination
preserves SPD).

Available backends in \texttt{LinearSolverFactory}, controlled by
\texttt{SimulationConfig}:

\begin{longtable}[]{@{}
  >{\raggedright\arraybackslash}p{(\columnwidth - 4\tabcolsep) * \real{0.3333}}
  >{\raggedright\arraybackslash}p{(\columnwidth - 4\tabcolsep) * \real{0.3333}}
  >{\raggedright\arraybackslash}p{(\columnwidth - 4\tabcolsep) * \real{0.3333}}@{}}
\toprule\noalign{}
\begin{minipage}[b]{\linewidth}\raggedright
Setting
\end{minipage} & \begin{minipage}[b]{\linewidth}\raggedright
Backend
\end{minipage} & \begin{minipage}[b]{\linewidth}\raggedright
Use when
\end{minipage} \\
\midrule\noalign{}
\endhead
\bottomrule\noalign{}
\endlastfoot
\texttt{use\_petsc=1} (default) & PETSc KSP, options from env
\texttt{PETSC\_OPTIONS} & Production runs (project default) \\
\texttt{use\_hypre\_boomeramg=1} & MFEM CG + HypreBoomerAMG & Large
meshes, MPI ≥ 16 \\
\texttt{use\_hypre\_block\_jacobi=1} & MFEM CG + HypreSmoother
(l1-Jacobi) & Cheap PC, comparable to plain CG \\
(none of the above) & MFEM CG without preconditioner & Diagnostic /
small problems \\
\end{longtable}

Project-default \texttt{config/petsc\_asm.opts}:

\begin{verbatim}
-mono_ksp_type cg
-mono_ksp_max_it 500
-mono_ksp_rtol 1e-8
-mono_ksp_norm_type unpreconditioned
-mono_pc_type asm
-mono_pc_asm_overlap 0
-mono_sub_pc_type icc
-mono_sub_pc_factor_levels 0
\end{verbatim}

ICC(0) (Incomplete Cholesky, zero fill) exploits SPD structure; overlap
0 is the cheapest Schwarz variant. On the half-ellipsoid 124k-tet
problem this gives mean \textbf{5 KSP iter/step} vs \textbf{50
iter/step} for plain CG (10× reduction); see
\texttt{docs/petsc\_asm\_vs\_cg\_iter\_bench.md}.

Activation:

\begin{Shaded}
\begin{Highlighting}[]
\BuiltInTok{export} \VariableTok{PETSC\_OPTIONS}\OperatorTok{=}\StringTok{"}\VariableTok{$(}\FunctionTok{grep} \AttributeTok{{-}v} \StringTok{\textquotesingle{}\^{}\#\#\textbackslash{}|\^{}$\textquotesingle{}}\NormalTok{ config/petsc\_asm.opts }\KeywordTok{|} \FunctionTok{tr} \StringTok{\textquotesingle{}\textbackslash{}n\textquotesingle{}} \StringTok{\textquotesingle{} \textquotesingle{}}\VariableTok{)}\StringTok{"}
\ExtensionTok{./build/monodomain} \AttributeTok{{-}{-}config}\NormalTok{ config/half\_ellipsoid\_purkinje.options}
\end{Highlighting}
\end{Shaded}


\end{document}
